\documentclass[12pt]{article}
\usepackage[utf8]{inputenc}
\usepackage[T1]{fontenc}
\usepackage[margin=1in]{geometry}
\usepackage{setspace}
\usepackage{xcolor}
\usepackage[colorlinks=true, urlcolor=blue!60!black, linkcolor=blue!60!black]{hyperref}
\setlength{\parindent}{2em}
\setlength{\parskip}{1.25em}
\renewcommand{\baselinestretch}{1.0}
\begin{document}
\begin{center}
{\LARGE \textit{Tennessee Wine \& Spirits Retailers Ass'n v. Thomas}, 588 U.S. 504 (2019)\par}
\vspace{0.8em}
{\large Case Brief --- Four-Page Edition\par}
\end{center}
\vspace{1em}

\section{Memory Jogger}

Tennessee required retail liquor-store license applicants to have lived in the State for two years; the Court held 7--2 that the rule violated the dormant Commerce Clause and was not saved by \S~2 of the Twenty-first Amendment, which does not shield protectionist alcohol laws lacking evidence-backed health or safety justifications (\href{https://www.supremecourt.gov/opinions/18pdf/18-96_5i36.pdf}{Slip Opinion}).

\section{Detailed Case Facts}

Tennessee's three-tier system separates producers, wholesalers, and retail package stores. Its residency rules required (1) two years of residency for an initial retail license, (2) ten consecutive years to renew, and (3) for corporations, that 100 percent of officers, directors, and stockholders meet those requirements, which barred publicly traded companies. In 2012 the Tennessee Attorney General found the rules unconstitutional, and the Alcoholic Beverage Commission (TABC) stopped enforcing them, even after the legislature added a ``health, safety and welfare'' intent clause.

In 2016, two nonqualifying applicants sought licenses: Total Wine, a Tennessee LLC owned by Maryland residents, and a Memphis store owned by Doug and Mary Ketchum, who had moved from Utah for their daughter's health. When TABC staff recommended approval, the Tennessee Wine and Spirits Retailers Association, a trade group of incumbent stores, threatened to sue.

\section{Procedural History}

The TABC's director sought a declaratory judgment; after removal, the Middle District of Tennessee held the rules unconstitutional under \textit{Granholm v. Heald} (2005), and the licenses issued. The Association appealed, and a divided Sixth Circuit (Judge Sutton dissenting in part) struck down all three provisions. Certiorari was sought only on the two-year rule. The Court affirmed on June 26, 2019.

\section{Judicial Votes}

Affirmed, \textbf{7--2}.
\begin{itemize}
\item \textbf{Majority:} Justice Alito, joined by Chief Justice Roberts and Justices Ginsburg, Breyer, Sotomayor, Kagan, and Kavanaugh.
\item \textbf{Concurrences:} None.
\item \textbf{Dissent:} Justice Gorsuch, joined by Justice Thomas.
\end{itemize}

\section{Holding}

The two-year requirement facially discriminates against nonresidents in violation of the dormant Commerce Clause and is not saved by \S~2, which does not shield protectionist alcohol regulations lacking a predominant, evidence-backed public health or safety justification.

\section{Analysis of Opinions}

\textbf{Majority (Alito): the dormant Commerce Clause.} Curbing state protectionism was a central reason for the Convention, and the dormant Commerce Clause remains the main bulwark against trade barriers. Outside the alcohol context, the residency rule would plainly fail.

\textbf{Section 2's history.} Section 2 bars ``[t]he transportation or importation into any State \ldots{} for delivery or use therein of intoxicating liquors, in violation of the laws thereof.'' Read literally, it would override every other guarantee. Instead it constitutionalized the Wilson Act and Webb-Kenyon Act, which let States regulate imported alcohol on equal terms but never authorized protectionism. \textit{Bacchus}, \textit{Healy}, and \textit{Granholm} confirm that protectionism is not a legitimate \S~2 interest.

\textbf{No retailer carveout.} \textit{Granholm} called the basic three-tier model ``unquestionably legitimate,'' but did not immunize every discriminatory add-on, and residency rules are not an ``essential feature'' of it.

\textbf{The test applied.} A law must ``be justified as a public health or safety measure or on some other legitimate nonprotectionist ground,'' shown by ``concrete evidence,'' not ``mere speculation.'' Each rationale failed: registered agents secure amenability to process; background checks screen fitness (the TABC did not investigate applicants during the two-year wait); inspections and revocation provide oversight; and ``responsible consumption'' was unsupported and underinclusive, since a Memphis resident could open a store in Bristol, Tennessee, but a resident of Bristol, Virginia, could not.

\textbf{Gorsuch (dissenting).} The dormant Commerce Clause is a ``peculiar'' implication that the people overrode for alcohol in \S~2. Residency laws date to New Hampshire in 1834, and 18 States enacted them after Repeal. \textit{Granholm} deliberately separated products from distribution, and the ``predominant effect'' standard gives lower courts no workable guidance.

\textbf{Later developments.} Circuits divide over in-state premises rules for wine shippers: some, as in \textit{Day v. Henry} (9th Cir. 2025), treat physical presence as an ``essential feature,'' while \textit{Block v. Canepa} (6th Cir. 2023) demands concrete evidence.

\section{Examples: Future Applications}

\begin{enumerate}
\item \textbf{Same-Side (Invalid): Property-Ownership Rule.} Licenses limited to three-year in-state property owners rest on ``mere speculation'' and are no essential feature of three tiers.
\item \textbf{Same-Side (Invalid): All-Resident Ownership.} Requiring every officer and shareholder to be a resident discriminates on its face; agents, checks, and inspections suffice.
\item \textbf{Opposite-Side (Likely Valid): Neutral Caps and Price Floors.} Evenhanded density caps, uniform price floors, and storefront rules backed by health studies fall within \S~2's leeway.
\item \textbf{Opposite-Side (Likely Valid): State Monopoly.} A state-run retail monopoly excludes all private sellers equally, so it triggers no discrimination inquiry.
\item \textbf{Fence-Sitter: Direct Shipping.} In-state retailers may ship to consumers but out-of-state retailers may not without an in-state store; the circuits are divided.
\end{enumerate}

\section{Critique}

\textbf{Doctrinal indeterminacy.} The opinion creates a hybrid standard, looser than ordinary dormant Commerce Clause scrutiny but stricter than rational-basis review, without explaining how to weigh mixed-motive laws. As Gorsuch noted, limiting entry and raising prices can reduce consumption, and the majority never explains when higher prices are legitimate temperance rather than protectionism.

\textbf{Text versus structure.} Gorsuch and Thomas prioritized \S~2's enacted text and its 1933 public meaning over an implied doctrine; Alito harmonized \S~2 with the Constitution's anti-Balkanization structure.

\section{Key Quotations}

\begin{quote}
``Section 2 gives the States regulatory authority that they would not otherwise enjoy, but as we pointed out in \textit{Granholm}, `mere speculation' or `unsupported assertions' are insufficient to sustain a law that would otherwise violate the Commerce Clause. \ldots{} Where the predominant effect of a law is protectionism, not the protection of public health or safety, it is not shielded by \S~2.''
\end{quote}
\noindent --- Justice Alito, for the Court, 588 U.S. at 539--540.

\begin{quote}
``Section 2 was adopted as part of the scheme that ended prohibition on the national level. It gives each State leeway in choosing the alcohol-related public health and safety measures that its citizens find desirable. But \S~2 is not a license to impose all manner of protectionist restrictions on commerce in alcoholic beverages.''
\end{quote}
\noindent --- Justice Alito, for the Court, slip op. at 2.

\begin{quote}
``At issue in the present case is not the basic three-tiered model of separating producers, wholesalers, and retailers, but the durational-residency requirement that Tennessee has chosen to impose on new applicants for liquor store licenses. Such a requirement is not an essential feature of a three-tiered scheme.''
\end{quote}
\noindent --- Justice Alito, for the Court, slip op. at 28.

\end{document}
